\begin{claim}{F.13}{Proved}{The exact homology kernel and the retraction obstruction}
Retain the specified injection
$\widehat\phi$ of F.12.2.  Write $e_0,e_\infty,e_j^+,e_j^-$
for the twelve positive source cusp classes, with the
single relation
\[
 e_0+e_\infty+\sum_{j=0}^4(e_j^++e_j^-)=0.
                                                               \tag{F.13.1}
\]
For $A=\Z$, $\widehat\Z$, or $\Z_p$ at any prime $p$,
the actual first-homology map induced by $\phi$ has exact
kernel
\[
 \ker\phi_{*,1}
   =A\,k,\qquad
 k=e_0-e_\infty
   =2e_0+\sum_{j=0}^4(e_j^++e_j^-).
                                                               \tag{F.13.2}
\]
Here the second expression uses the genuine basis
omitting $e_\infty$.  In particular $k$ is primitive
over $\Z$ and has nonzero reduction modulo every prime.

The discrete map, its full profinite completion, and its
maps to maximal pro-$p$ completions have no homomorphic
left inverse on the stated domains, continuous where
profinite groups are involved.  The same obstruction
holds after every further target-forgetting map.
No assertion about a finite-index target subgroup, the
injectivity of a forgotten map, or the full equalizer
F.12.5 is made.
\end{claim}
\begin{proof}
To prove the kernel equality, write an arbitrary class
in that basis as
\[
 z=c e_0+\sum_{j=0}^4(p_j e_j^++m_j e_j^-).
\]
The twenty target winding characters are the actual
integral basis of F.12.14.  Their pullbacks from
F.12.3 and F.12.15 give
\[
 \begin{aligned}
 L_{j+1}(\phi_*z)&=2(p_j-m_j),\\
 M_{j+1}(\phi_*z)&=c-2m_j,\\
 N_{i+1,j+1}(\phi_*z)
   &=c+p_k+m_k-2m_i-2m_j,
      &2k&\equiv i+j\pmod5\quad(i\ne j).
 \end{aligned}                                                \tag{F.13.3}
\]
Multiplication by two is injective in each of the
specified coefficient rings.  Vanishing of the first
two rows therefore gives $p_j=m_j=t$ for all $j$ and
$c=2t$.  The third row then vanishes identically.
This proves (F.13.2), without identifying an actual
group kernel with its abelian image.

For completeness, at a specified primary coordinate
$p$, the vector $k$ has ten entries equal to one in
the displayed $\Z_p$-basis.  It is consequently
nonzero over $\Z_p$ and modulo $p$.  In the full
$\widehat\Z$-module, multiplying $k$ by the idempotent
whose $p$-coordinate is one and whose other primary
coordinates are zero gives a nonzero kernel vector
supported exactly at that primary coordinate.

There is no homomorphic left inverse of the discrete
map $\phi_*:\pi_1(Y)\to P\Gamma_{0,8}$ and no continuous
homomorphic left inverse of its full profinite map.
Indeed such a left inverse would induce a left inverse
on abelianizations, contradicting (F.13.2).
More precisely, for every prime $p$ there is no
continuous homomorphism
\[
 r_p:\widehat{P\Gamma}_{0,8}
                  \longrightarrow\pi_1(Y)^{(p)}
 \quad\hbox{with}\quad
 r_p\widehat\phi=\pi_p,
                                                               \tag{F.13.4}
\]
where $\pi_p$ is the actual maximal pro-$p$ quotient
map.  Apply abelianization modulo $p$ to this equality:
its left side kills $k\bmod p$, whereas its right side
is the identity on the source's $\F_p$-homology.
The same argument rules out a left inverse of the
induced map between maximal pro-$p$ completions.

Every homomorphism obtained by first applying
$\widehat\phi$ and then forgetting any target labels
also kills $k$ on first homology, by functoriality.
Thus none of those composites has a homomorphic
left inverse.  This does not decide which such
composites are injective, nor does it decide the
equalizer F.12.5.  The obstruction is stated for
these displayed domains; passing to a finite-index
target subgroup would change its first-homology
map and is not covered by this argument.

The separate exact certificate
\path{code/part_a/common_reflection_retraction.py}
constructs all twenty divisor rows from (F.13.3),
checks the source puncture-sum relation and the
primitive kernel vector, computes rank ten over
$\Q$, and checks the explicit nonzero kernel vector
in the listed finite primary reductions.  Its
finite reductions verify the displayed formulas;
the uniform statement for every prime follows
from the coefficient-one argument above.
\end{proof}
